% !PRINTED 571
\weberpage{571}
%<<CONT
halten, und wenn wir also die Variablen von $\Phi$ mit $t, u_1, u_2,
\ldots$ bezeichnen, werden wir auch setzen
\weq{5}{\Phi(t) = \Phi(t, u_1, u_2, \ldots),}
oder kürzer $\Phi(t, u)$, und diese Function enthält ganze rationale
Zahlencoëfficienten. Wenn wir die Function $\Phi(t)$ in die
$p^{\text{te}}$ Potenz erheben, und die Formel \S.\,150, 4. anwenden,
so folgt aus der Congruenz (3) eine neue Congruenz
\weq{6}{\Phi(\tau^p, u_1^p, u_2^p, \ldots) \equiv 0\ (\mathrm{mod}\ \mathfrak p).}

Ebenso ist aber auch
\weq{7}{\tau^p \equiv \omega_1^p t_1^p + \omega_2^p t_2^p + \cdots + \omega_n^p t_n^p.}

Wenn wir dies in die Congruenz (6) substituiren, so entsteht eine
durch $\mathfrak p$ theilbare Function, in der die Variablen
$u_1, u_2, \ldots$, unter denen ja die $t_1, t_2, \ldots, t_n$ mit
enthalten sind, nur in der $p^{\text{ten}}$ Potenz vorkommen. Die
Congruenz muss also richtig bleiben, wenn wir die $u_1^p, u_2^p,
\ldots$ und also auch die $t_1^p, t_2^p, \ldots, t_n^p$ durch
unabhängige Variable $u_1, u_2, \ldots t_1, t_2, \ldots, t_n$ ersetzen
(\S.\,143).

Setzen wir demnach
\weq{8}{\tau_1 = \omega_1^p t_1 + \omega_2^p t_2 + \cdots + \omega_n^p t_n,}
so ergiebt sich aus (6)
\[
\Phi(\tau_1, u_1, u_2, \ldots) \equiv 0\ (\mathrm{mod}\ \mathfrak p),
\]
und folglich ist $\tau_1$ auch eine Wurzel der Congruenz (4).

Dieses nämliche Verfahren lässt sich wiederholt anwenden, und wir
finden, dass auch
\[
\tau_2 = \omega_1^{p^2} t_1 + \omega_2^{p^2} t_2 + \cdots + \omega_p^{p^2} t_n % [sic: Weber prints \omega_p here; (9) has \omega_n.]
\]
Wurzel der Congruenz (4) ist, u.\,s.\,f.

Wenn wir also ein System von
Formen $\tau_r$ definiren durch
\weq{9}{\tau_r = \omega_1^{p^r} t_1 + \omega_2^{p^r} t_2 + \cdots + \omega_n^{p^r} t_n}
für beliebige positive Exponenten $r$, so sind alle diese Grössen
$\tau_r$ zugleich Wurzeln der Congruenz (4). Es ist noch die Frage zu
beantworten, wie viele von diesen Formen $\tau_r$ von einander
verschieden sind.

Nach \S.\,150, (4) ist sicher
\[
\tau_r \equiv \tau_{r'}\ (\mathrm{mod}\ \mathfrak p),
\]
wenn
\[
r \equiv r'\ (\mathrm{mod}\ f)
\]
ist, und folglich giebt es unter den $\tau_r$ gewiss nicht mehr als
$f$ nach dem Modul $\mathfrak p$ verschiedene
\weq{10}{\tau, \tau_1, \tau_2, \ldots, \tau_{f-1}.}
